\FINALFIX{\newpage}
\subsection{2D Couette Flow}
\label{SEC:Testcases:CouetteFlow}

\begin{figure}[htb]
  \begin{centering}
    \insertimage[width=0.3\textwidth]{tikz/couette_schematic}
   %\input{Figures/test_cases/couette_schematic}
    \caption[Couette flow test case.]{Schematic of Couette flow test case.}
    \label{FIG:CouetteFlowSchematic}
  \end{centering}
\end{figure}

Couette flow is a classic problem consisting of viscous flow in a
two-dimensional channel between two parallel walls, with one wall
moving at a constant velocity and the other wall stationary,
as shown in \fref{FIG:CouetteFlowSchematic}.
%
For the specific case described here, both walls are treated as
isothermal with a constant temperature of 300 Kelvin.
%
The top wall has a velocity of 70 m/s.

%\begin{nmlblocklst}[
,label={LST:testcase:couette}
,caption={Sample input file for Couette flow test case.}
]
&nml_version
  Namelist_Version          =  2
/
&nml_grid
  Grid_Format               =  3  ! Plot3D format
  Grid_File                 =  '../grid_files/couette.grd'
  Plot3d_BCSfile            =  '../grid_files/couette.bcs.grd'
/
&nml_reference
  Mach                      =       0.201569299302318  ! U = 70 m/s
  Reynolds                  =     200.0
  Gamma                     =       1.4
  Gas_Constant              =     287.142857142857     ! J/(kg*K)
  Temperature               =     300.0  ! K
  Pressure                  =  100000.0  ! Pa
  Suth_Muref                =       0.406301824212273  ! Pa*s
  Suth_Tref                 =     300.0  ! K
  Suth_Sref                 =       0.0  ! K
/
&nml_physics
  Governing_Equations       =  2  ! Navier-Stokes
/
&nml_initialization
  Initialization            =  1  ! Init to ref conditions
/
&nml_timestep
  Runge_Kutta_Scheme        =  3
  Num_Timesteps             =  1000000
  Final_Time                =    100.0
  Timestep_Type             =  1  ! Global CFL
  CFL                       =  0.1
  CFL_Ramping_Beg           =  0.1
  CFL_Ramping_End           =  0.4
  CFL_Ramping_Iterations    =  1000
  Convergence_Order_Abs_Res = -10.0
  Convergence_Order_Max_Res = -10.0
/
&nml_numerics
  All_Orders                =  2  ! P2 solution, 3rd order accuracy
  All_Points                =  1  ! Legendre-Gauss points
/
&nml_output
  Restart_Interval          =  10000
  Display_Interval          =    100
  Residual_Interval         =    100
  Output_Interval           =  10000
/
&nml_bc
  BC_Input(1)%Name                      =  'TOP_WALL'
  BC_Input(1)%Set_BC_Value              =   53         ! Isothermal Wall
  BC_Input(1)%VelocityX                 =   70.0       ! m/s
  BC_Input(1)%Wall_Temp                 =  300.0       ! K
  !
  BC_Input(2)%Name                      =  'BOT_WALL'
  BC_Input(2)%Set_BC_Value              =   53         ! Isothermal Wall
  BC_Input(2)%Wall_Temp                 =  300.0       ! K
  !
  BC_Input(3)%Name                      =  'INFLOW'
  BC_Input(3)%Set_BC_Value              =   81         ! Periodic
  BC_Input(3)%Periodic%Matching_BC_Name =  'OUTFLOW'
  BC_Input(3)%Periodic%Matching_BC_Idx  =    4
  BC_Input(3)%Periodic%Translation(1)   =    2.0
  !
  BC_Input(4)%Name                      =  'OUTFLOW'
  BC_Input(4)%Set_BC_Value              =   81         ! Periodic
  BC_Input(4)%Periodic%Matching_BC_Name =  'INFLOW'
  BC_Input(4)%Periodic%Matching_BC_Idx  =    3
  BC_Input(4)%Periodic%Translation(1)   =   -2.0
/
\end{nmlblocklst}



